\documentclass[9pt,landscape]{extarticle}
\usepackage[letterpaper,margin=0.45in]{geometry}
\usepackage{multicol}
\usepackage{amsmath,amssymb}
\usepackage{multirow}
\usepackage{listings}
\usepackage{booktabs}
\usepackage[table]{xcolor}
\usepackage{titlesec}
\usepackage{enumitem}

\definecolor{kw}{RGB}{0,92,197}
\definecolor{accent}{RGB}{160,40,140}
\definecolor{accent2}{RGB}{90,70,160}
\definecolor{cmt}{RGB}{110,110,110}
\definecolor{rowgray}{gray}{0.9}

\lstdefinestyle{py}{
  language=Python,
  basicstyle=\ttfamily\footnotesize,
  keywordstyle=\color{kw},
  morekeywords={None,True,False,lambda},
  commentstyle=\color{cmt}\itshape,
  stringstyle=\color{accent},
  showstringspaces=false,
  columns=fullflexible,
  keepspaces=true,
  breaklines=true,
  breakatwhitespace=true,
  aboveskip=3pt,
  belowskip=3pt,
  xleftmargin=0pt,
}
\lstset{style=py}

\titleformat{\section}{\large\bfseries}{}{0pt}{}
\titlespacing*{\section}{0pt}{6pt}{3pt}
\titleformat{\subsection}{\normalsize\bfseries}{}{0pt}{}
\titlespacing*{\subsection}{0pt}{4pt}{2pt}

\setlength{\parindent}{0pt}
\setlength{\parskip}{2pt}
\setlength{\columnsep}{16pt}
\pagestyle{empty}

\newcommand{\hl}[1]{\textcolor{accent}{#1}}
\newcommand{\hll}[1]{\textcolor{accent2}{#1}}
\newcommand{\kw}[1]{\textcolor{kw}{\texttt{#1}}}

\begin{document}
\begin{multicols*}{3}

\fbox{\parbox{0.95\linewidth}{\centering
{\Large\bfseries CS1101S Cheatsheet}\\
for midterms AY26-27\\
\textcolor{accent2}{adapted from m.\ zaidan}}}

\section*{Recursion \& Iteration}
Recursion: \hl{Increasing} deferred operations\\
Iteration: \textcolor{kw}{Constant} deferred operations\\
(rule of thumb: if the final call is the application call, it is usually \emph{iterative})

\section*{Definitions of List \& Tree}
\begin{lstlisting}
pair = lambda x, y: (lambda f: f(x, y))
head = lambda p: p(lambda x, y: x)
tail = lambda p: p(lambda x, y: y)
\end{lstlisting}

A linked list is either \kw{None} or \hl{a pair whose tail is a linked list.}\\
A tree is either \kw{None} or \hl{head is an element, tail is a tree} or \hll{head is a tree, tail is a tree}

\subsection*{Source built-ins, in Python}
\begin{lstlisting}

def is_llist(xs):
    return (True if is_none(xs)
            else is_llist(tail(xs)) if is_pair(xs)
            else False)

def length(xs):
    return 0 if is_none(xs) else 1 + length(tail(xs))

def append(xs, ys):
    return (ys if is_none(xs)
            else pair(head(xs), append(tail(xs), ys)))

def remove(x, xs):          # first occurrence only
    return (None if is_none(xs)
            else tail(xs) if head(xs) == x
            else pair(head(xs), remove(x, tail(xs))))
\end{lstlisting}

\section*{List Abstractions}
\begin{lstlisting}
def map(f, xs):
    return (None if is_none(xs)
            else pair(f(head(xs)), map(f, tail(xs))))
\end{lstlisting}
Recursive process; \textbf{time: O(n), space: O(n)}, where n is the length of xs.

\begin{lstlisting}
def filter(pred, xs):
    return (None if is_none(xs)
            else pair(head(xs), filter(pred, tail(xs)))
                 if pred(head(xs))
            else filter(pred, tail(xs)))
\end{lstlisting}
Recursive process; \textbf{time: O(n), space: O(n)}, where n is the length of xs.

\begin{lstlisting}
def reduce(op, initial, xs):
    return (initial if is_none(xs)
            else op(head(xs),
                    reduce(op, initial, tail(xs))))

reduce(lambda curr, wish: ..., initial, xs)
\end{lstlisting}
Recursive process; \textbf{time: O(n), space: O(n)}, where n is the length of xs assuming op takes constant time.

\section*{Tree Abstractions}
\begin{lstlisting}
def reduce_tree(f, op, initial, tree):
    return reduce(lambda x, ys:
            op(reduce_tree(f, op, initial, x), ys)
            if is_llist(x)
            else op(f(x), ys),
        initial, tree)

def map_tree(f, tree):
    return map(lambda sub_tree:
            f(sub_tree) if not is_llist(sub_tree)
            else map_tree(f, sub_tree),
        tree)

def flatten_tree(xs):
    def h(xs, prev):
        return (prev if is_none(xs)
                else append(flatten_tree(xs), prev)
                     if is_llist(xs)
                else pair(xs, prev))
    return reduce(h, None, xs)
\end{lstlisting}

\subsection*{Binary Search Trees}
\begin{lstlisting}
def make_empty_tree(): return None
def is_empty_tree(t): return is_none(t)
def make_tree(e, l, r): return llist(e, l, r)
def entry(t): return head(t)
def left_branch(t): return head(tail(t))
def right_branch(t): return head(tail(tail(t)))

def insert(bst, item):
    if is_empty_tree(bst):
        return make_tree(item, make_empty_tree(),
                               make_empty_tree())
    elif item < entry(bst):
        return make_tree(entry(bst),
                         insert(left_branch(bst), item),
                         right_branch(bst))
    elif item > entry(bst):
        return make_tree(entry(bst),
                         left_branch(bst),
                         insert(right_branch(bst), item))
    else:
        return bst

def find(bst, name):
    return (False if is_empty_tree(bst)
            else True if name == entry(bst)
            else find(left_branch(bst), name)
                 if name < entry(bst)
            else find(right_branch(bst), name))
\end{lstlisting}

\section*{Permutations \& Combinations}
\begin{lstlisting}
def permutations(s):
    return (llist(None) if is_none(s)
        else reduce(append, None,
            map(lambda x: map(lambda p: pair(x, p),
                        permutations(remove(x, s))),
                s)))

def subsets(s):
    return reduce(
        lambda x, s1: append(s1,
            map(lambda ss: pair(x, ss), s1)),
        llist(None),
        s)

def choose(n, r):
    if n < 0 or r < 0:
        return 0
    elif r == 0:
        return 1
    else:
        to_use = choose(n - 1, r - 1)
        not_to_use = choose(n - 1, r)
        return to_use + not_to_use

def combinations(xs, r):
    if (r != 0 and is_none(xs)) or r < 0:
        return None
    elif r == 0:
        return llist(None)
    else:
        no = combinations(tail(xs), r)
        yes = combinations(tail(xs), r - 1)
        yes_item = map(lambda x: pair(head(xs), x),
                       yes)
        return append(no, yes_item)

def makeup_amount(x, coins):
    if x == 0:
        return llist(None)
    elif x < 0 or is_none(coins):
        return None
    else:
        combi_A = makeup_amount(x, tail(coins))
        combi_B = makeup_amount(x - head(coins),
                                tail(coins))
        combi_C = map(lambda x: pair(head(coins), x),
                      combi_B)
        return append(combi_A, combi_C)
\end{lstlisting}

\section*{Rule of Thumb for Abstractions}
{\small
\begin{tabular}{@{}ll@{}}
Is input a llist? & if not, don't use \\
Is length(output) = length(input)? & use \emph{map} \\
Are items in output = items in input? & use \emph{filter} \\
Else & use \emph{reduce}
\end{tabular}}

\section*{Useful Math Functions}
\begin{lstlisting}
import math
math.pow(base, exponent)   # or base ** exponent
round(x)       # banker's rounding: round(2.5) == 2
math.floor(x)  # a // b is floor division
math.ceil(x)
math.sqrt(x)
\end{lstlisting}

\newpage


\section*{Orders of Growth}
for $r(n)$,

\textbf{Big O} $O(g(n))$: if there is a positive constant $k$ such that $r(n) \le k \cdot g(n)$ for any sufficiently large value of $n$\\
\textbf{Big Theta} $\Theta(g(n))$: if there are positive constants $k_1$ and $k_2$ and a number $n_0$ such that $k_1 \cdot g(n) \le r(n) \le k_2 \cdot g(n)$ for any $n > n_0$.\\
\textbf{Big Omega} $\Omega(g(n))$: if there is a positive constant $k$ such that $k \cdot g(n) \le r(n)$ for any sufficiently large value of $n$\\
\textbf{Order (ascending):} $1, \log n, n, n\log n, n^2, n^3, 2^n, 3^n, n^n$

Ignore constants, lower order terms. $O(2n) = O(3n) = O(n)$\\
For a sum, take the larger term. $O(n) + O(n^2) = O(n^2)$\\
For a product, multiply the two terms. $O(n) \times O(n) = O(n^2)$

\section*{Recurrence Relations}
{\small
\begin{tabular}{@{}cc@{}}
\toprule
$T(n)$ & Solution \\
\midrule
$O(1) + T(n-1)$ & \\
$O(1) + 2T(n/2)$ & $O(n)$ \\
$O(n) + T(n/2)$ & \\
\rowcolor{rowgray} $O(1) + T(n/2)$ & $O(\log n)$ \\
$O(\log n) + T(n-1)$ & \\
$O(n) + 2T(n/2)$ & \multirow{-2}{*}{$O(n \log n)$} \\
\rowcolor{rowgray} $O(n) + T(n-1)$ & $O(n^2)$ \\
$O(n^k) + T(n-1)$ & $O(n^{k+1})$ \\
\rowcolor{rowgray} $O(n) + 2T(n-1)$ & $O(2^n)$ \\
\bottomrule
\end{tabular}}

\section*{Master Theorem}
For $T(n) = a\,T(n/b) + f(n)$, with $a \ge 1$, $b > 1$:\\
$a$ = no.\ of subproblems, $n/b$ = size of each, $f(n)$ = work to split and combine.

\textbf{Shortcut} when $f(n) = \Theta(n^d)$: compare $a$ with $b^d$.\\
$a < b^d \Rightarrow \Theta(n^d)$ \hfill root wins\\
$a = b^d \Rightarrow \Theta(n^d \log n)$ \hfill tie\\
$a > b^d \Rightarrow \Theta(n^{\log_b a})$ \hfill leaves win

{\small
\begin{tabular}{@{}lccl@{}}
\toprule
Recurrence & $a,b,d$ & Case & Result \\
\midrule
Binary search: $T(n/2) + O(1)$ & $1,2,0$ & tie & $\Theta(\log n)$ \\
\rowcolor{rowgray} Merge sort: $2T(n/2) + O(n)$ & $2,2,1$ & tie & $\Theta(n\log n)$ \\
Tree traversal: $2T(n/2) + O(1)$ & $2,2,0$ & leaves & $\Theta(n)$ \\
\rowcolor{rowgray} $T(n/2) + O(n)$ & $1,2,1$ & root & $\Theta(n)$ \\
\bottomrule
\end{tabular}}

Only for splitting by a \emph{factor} ($n/b$). For $T(n-1)$ recurrences, use the table above.

\section*{Sorting/Search Algorithms}
{\small
\begin{tabular}{@{}lccc@{}}
\toprule
 & \multicolumn{2}{c}{\textbf{Time}} & \textbf{Space} \\
 & best & worst & \\
\midrule
Binary Search & $\Theta(\log n)$ & $O(\log n)$ & $O(1)$ \\
\rowcolor{rowgray} Selection Sort & $\Theta(n^2)$ & $O(n^2)$ & $O(n)$ \\
Insertion Sort & $\Theta(n)$ & $O(n^2)$ & $O(n)$ \\
\rowcolor{rowgray} Merge Sort & $\Theta(n\log n)$ & $O(n\log n)$ & $O(n\log n)$ \\
Quick Sort & $\Theta(n\log n)$ & $O(n^2)$ & $O(n^2)$ \\
\bottomrule
\end{tabular}}

\subsection*{Selection Sort}
\begin{lstlisting}
def smallest(xs):
    return reduce(lambda x, y: x if x < y else y,
                      head(xs), tail(xs))

def selection_sort(xs):
    if is_none(xs):
        return xs
    else:
        x = smallest(xs)
        return pair(x, selection_sort(remove(x, xs)))
\end{lstlisting}

\subsection*{Merge Sort}
\begin{lstlisting}
def merge(xs, ys):
    if is_none(xs):
        return ys
    elif is_none(ys):
        return xs
    else:
        x = head(xs)
        y = head(ys)
        return (pair(x, merge(tail(xs), ys)) if x < y
                else pair(y, merge(xs, tail(ys))))

middle = lambda n: n // 2     # math.floor(n / 2)

def take(xs, n):
    return (None if n == 0
            else pair(head(xs), take(tail(xs), n - 1)))

def drop(xs, n):
    return xs if n == 0 else drop(tail(xs), n - 1)

def merge_sort(xs):
    if is_none(xs) or is_none(tail(xs)):
        return xs
    else:
        m = middle(length(xs))
        return merge(merge_sort(take(xs, m)),
                     merge_sort(drop(xs, m)))
\end{lstlisting}

\subsection*{Quick Sort}
\begin{lstlisting}
def partition(xs, p):
    return pair(filter(lambda x: x <= p, xs),
                filter(lambda x: x > p, xs))

def quicksort(xs):
    if is_none(xs):
        return None
    elif is_none(tail(xs)):
        return xs
    else:
        pivot = head(xs)
        ptn = partition(tail(xs), pivot)
        return reduce(append, None,
            llist(quicksort(head(ptn)), llist(pivot),
                 quicksort(tail(ptn))))
\end{lstlisting}

\section*{Continuation-Passing Style (CPS)}
Passing the deferred operation as a function in an extra argument. Any recursive function can be converted this way.

\begin{lstlisting}
def fac_cps(n, c):
    return (c(1) if n == 1
            else fac(n - 1, lambda x: c(n * x)))

def append_cps(xs, ys, c):
    return (c(ys) if is_none(xs)
            else append_cps(tail(xs), ys,
                lambda res: c(pair(head(xs), res))))

def reduce_cps(op, initial, xs, c):
    return (c(initial) if is_none(xs)
            else reduce_cps(op, initial, tail(xs),
                lambda res: c(op(head(xs), res))))

\end{lstlisting}

\textbf{Trace:} \kw{fac(3, I)} $\to$ \kw{fac(2, lambda r: I(3*r))} $\to$ \kw{fac(1, lambda r: I(3*(2*r)))} $\to$ \kw{fac(0, \dots)} $\to$ \kw{k(1)} $=$ \kw{3*(2*(1*1))} $=6$

\small{I = lambda x: x}

Iterative process, but \textbf{space is O(n).}

\end{multicols*}
\end{document}